\mxunit{terminale-sources}{Sources et périmètre}
\textbf{Sources fournies.} Barbazo, \emph{Terminale spécialité mathématiques}, Hachette Livre, 2020, ISBN 978-2-01-786619-0 ; progression annuelle de terminale spécialité 2026--2027 (15 pages) ; programme de spécialité de terminale fourni (nouvelle annexe) ; éphéméride 2026--2027 de l'établissement. Le manuel fournit les exercices précisément référencés dans les chapitres 1 à 11 ; la progression fournit l'ordre et les identifiants de séances.
\textbf{Agenda.} Événements « Maths spé », du 4 septembre 2026 au 2 juillet 2027, lus le 5 septembre 2026. Les descriptions sont prises comme intentions générales ; les contenus précis et les prérequis suivent le PDF de progression.
\textbf{Distinction entre les programmes.} Le PDF fourni présente la nouvelle organisation (vocabulaire et logique avant les domaines) et mentionne notamment les plans perpendiculaires. Il ne doit pas être confondu avec l'annexe de 2019 applicable à 2026--2027. Le contrôle du périmètre a été repris sur cette dernière ; les plans perpendiculaires sont signalés comme complément, sans évaluation dédiée.\par
\href{https://eduscol.education.gouv.fr/sites/default/files/document/spe246annexe1158907pdf-84159.pdf}{Programme de spécialité de terminale : annexe de 2019}\par
\href{https://www.education.gouv.fr/bo/2026/Hebdo14/MENE2602919A}{Arrêté du 26 février 2026 : application à la rentrée 2027--2028}\par
\textbf{Vérification officielle.} Page Éduscol « Programmes et ressources en mathématiques -- voie GT », consultée les 5 et 6 septembre 2026 : programme de terminale de 2019 pour cette année et nouveau programme à partir de 2027--2028.\par
\href{https://eduscol.education.gouv.fr/5817/programmes-et-ressources-en-mathematiques-voie-gt}{Éduscol : programmes et ressources en mathématiques}
\textbf{Référence d'épreuve récente.} Sujet ministériel de spécialité mathématiques, Métropole, 17 juin 2025 : durée 4 h et quatre exercices obligatoires. Utilisé comme repère de rédaction et comme travail de référence : exercice 1, questions 5 et 6 ; exercice 3. Leurs énoncés sont à lire dans l'annale, leurs corrigés figurent dans l'édition professeur.\par
\href{https://www.education.gouv.fr/sites/default/files/2025-06/baccalaur-at-g-n-ral-2025-math-matiques-440856.pdf}{Ministère : sujet de spécialité du 17 juin 2025}
\textbf{Calendrier national.} La page ministérielle de la session 2027 a été consultée ; elle ne remplace pas les instructions particulières aux centres étrangers ni la convocation individuelle. Le livret ne fixe donc pas de date d'examen local sur cette seule base.\par
\href{https://www.education.gouv.fr/reussir-au-lycee/baccalaureat-brevet-cap-le-calendrier-2027-341384}{Ministère : calendrier national de la session 2027}
\textbf{Repères de rédaction.} Les énoncés ont été repris en explicitant les objets, domaines, hypothèses et résultats attendus, et en séparant conjecture et démonstration. Exemples de rédaction consultés :\par
\href{https://www.maths-et-tiques.fr/telech/20SuitesTS2.pdf}{Maths et tiques, Yvan Monka : limites de suites}\par
\href{https://www.lyceedadultes.fr/sitepedagogique/documents/math/mathTermSpe/ctrle_2020_2021/13_devoir_01_03_2021.pdf}{Lycée d'adultes, Paul Milan : devoir du 1er mars 2021}\par
\textbf{Créations du livret.} Notes, activités, exercices, DM, sujets zéro et évaluations sont rédigés pour ce document. Les modèles contextualisés sont des situations pédagogiques, pas des données empiriques. Le calendrier de distribution et les choix de séance sont explicités pour permettre leur adaptation au rythme réel de la classe.
\end{document}


